\section{Discussion}
\label{sec:discussion}

\subsection{Why Fan-Out Points the Wrong Way}

The direction of propagation is the result that needs the fewest qualifications. In
\val{descendant_affected_runs_zero} of \val{campaign_runs_total} runs no descendant of the
faulted service degraded, and which services degraded hardly depended on how the fault was
delivered: for \val{fault_invariant_services_all} of the
\val{fault_invariant_services_all_total} services the per-run counts were identical whether the
service was killed or slowed down. The most plausible reading is that a caller suffers while it
waits on a dependency that has died or slowed, whereas a callee only receives fewer requests and
has nothing to wait for. The measurements fit that reading without testing it, because the
degradation rule asks only whether a service's own latency or error rate crossed a threshold and
never observes waiting directly.

If the reading is right, a fan-out ranking scores the wrong end of each edge. The composition
service of Social Network is the clearest case: it calls \val{composepost_out_degree} services
and the fan-out metric ranks it first, yet in each of its \val{composepost_runs_total} runs
exactly one caller degraded, and that caller was always the gateway.
%[id:composepost_runs_only_gateway_ancestor]
Across all services, the metric's correlation with degraded callers was negative in every cell
and separable from zero in none (Section~\ref{sec:res-h2}), and no weighting on the
preregistered grid changed that (Section~\ref{sec:res-weights}). The data give the fan-out
premise no support; they cannot show a negative relationship.

\subsection{A Predictor That Sits Close to Its Outcome}

The caller count correlates with degraded callers partly by construction, since the outcome
counts a subset of the predictor, and the two are equal whenever every ancestor of a faulted
service degrades. That held for every Hotel Reservation service and for ten of the eleven
Social Network services. %[id:saturation_hr_kill_service] %[id:saturation_sn_kill_service]
A post hoc breakdown sharpens the point. In all \val{gw_runs_all} runs the gateway was among the
degraded ancestors, %[id:gw_runs_gateway_degraded_all]
and in both applications every analysable service has the gateway among its graph ancestors.
%[id:gw_sn_services_gateway_ancestor] %[id:gw_hr_services_gateway_ancestor]
Every measured ancestor count therefore includes one unit that is the same for every service,
and what remains to be explained comes from intermediate callers. In Social Network only
user-service departs from its graph count in a consistent way. Its two ancestors that degraded
sit on request paths the workload exercises, and its other two showed no degradation in any run,
an explanation formed after seeing the data and not tested here.

With user-service withheld, the Social Network correlation becomes \val{loo_userservice_rho}.
The result can fairly be read as resting on one service's worth of independent information. One
lesson may reach beyond this study: a structural predictor tested against a structurally defined
outcome needs a saturation measure reported next to the correlation, without which the
correlation cannot be interpreted.

\subsection{What Users Would See}
\label{sec:disc-users}

The preregistered outcome counts degraded services and says nothing about how many requests
fail. To come closer to what a user experiences, the gateway's end-to-end error rate under kill
faults was examined after the fact. This analysis is exploratory and descriptive.

Gateway error followed the request mix. For each faulted service, the Locust task weights of the
request types whose traces pass through it were summed and divided by the total weight. The rank
correlation between that share and the median gateway error was
\val{uf_sn_rho_err_weight_share} on Social Network and \val{uf_hr_rho_err_weight_share} on
Hotel Reservation, reported without p-values. In Hotel Reservation, killing reservation, which
sits on request types carrying \val{uf_hr_reservation_weight_share} of the weight, gave a median
gateway error rate of \val{uf_hr_reservation_err_median}, while killing recommendation, with a
share of \val{uf_hr_recommendation_weight_share}, gave \val{uf_hr_recommendation_err_median}.
The mix does not explain everything. The Social Network services reached only through the
compose request all share a weight of \val{uf_sn_compose_post_service_weight_share}, and their
median error rates still range from \val{uf_sn_unique_id_service_err_median} to
\val{uf_sn_compose_post_service_err_median}.

The hybrid metric's exploratory correlation with gateway error on Social Network was
\val{uf_test_sn_hybrid_rho}, not significant after correction within its family (adjusted
$p = \val{uf_test_sn_hybrid_bh_p}$). Splitting the metric into its weighted terms shows where
the association comes from: the in-degree term alone has a rank correlation of
\val{uf_sn_rho_err_term_in} with gateway error, the fan-out term \val{uf_sn_rho_err_term_fanout}
and the betweenness term \val{uf_sn_rho_err_term_btw}. The association follows the count of
direct callers more closely than the fan-out correction, so it lends that correction no support.

\subsection{Lessons about Measurement}

Two measurement choices mattered more than the plan anticipated. The first is where latency is
attributed. Charging each trace's end-to-end latency to every service on the trace makes services
on a shared path indistinguishable; in the pilot comparison of Section~\ref{sec:attribution} that
shortcut flagged five non-faulted services where per-service attribution flagged one.
%[fig:fig_q3_artifact]
The correction was in place before the campaign, so what an uncorrected campaign would have
concluded is unknown, but the pilot shows how it would have behaved.

The second is that per-service attribution is only as good as the instrumentation behind it.
Social Network's spans carry no span-kind tag, so its per-service latency drew on all of a
service's spans, \val{sn_spans_selected_client_pct}\% of them client spans. A post hoc replay on
server spans altered \val{replay_runs_changed} of the \val{replay_runs} run outcomes and none of
the service-level medians. %[id:replay_service_medians_changed]
The change was small, which was fortunate. A study that depends on span semantics should check
them before collecting data, and this one did not.

\subsection{Implications for Practitioners}

For a team that has to rank services before running any experiment, the data favour looking up
the graph at callers over looking down it at callees. That advice rests on the propagation
pattern and on the lack of support for fan-out; the preregistered test of the caller ranking
itself did not reach significance, so no claim is made that it predicts failure impact. In both
applications the gateway appeared in every caller count, and in the exploratory analysis the
user-visible error rate tracked how much traffic passed through the failed service. A ranking
that weights callers by the traffic they carry may therefore be worth testing as a new
preregistered hypothesis; it was not tested here. Larger applications would also lift the
ceiling that eleven services place on any correlation, and containment mechanisms such as retries
with back-off or circuit breakers, which were neither configured nor studied, could change how
far a failure travels.
